\documentclass{article}
\usepackage{../math_template}

\author{Jonathan Kasongo}
\title{Putnam and Beyond}

\begin{document}
\maketitle

\textit{Problems taken from the book by the same name as the title,
by Andreescu and Gelca.}

\begin{problem}{Chapter 2.2 --- Problem 147}
Let $P(x)$ be a polynomial of odd degree with real coefficients. Show
that $P(P(x)) = 0$ has at least as many real roots as the equation
$P(x) = 0$, counted without multiplicities.
\end{problem}

\begin{solution}{Solution}
It is well-known that any polynomial with odd degree has at least one real
root. Suppose $r_1 < r_2 < \cdots < r_k$ are the real roots of $P(x)$.
Consider the polynomials $P_i(x) = P(x) - r_i$, $(i=1,2,\ldots,k)$.
Noticing that $\deg P_i = \deg P$ we conclude that each $P_i(x)$ has
at least one real root. Suppose that $\alpha_i \in \mathbb{R}$ is one of
the real roots of $P_i$. Then
\[
\begin{aligned}
P(\alpha_i) &= r_i \\
P(P(\alpha_i)) &= P(r_i) = 0. \\
\end{aligned}
\]
So $\alpha_1, \alpha_2, \ldots, \alpha_k$ are roots of $P(P(x))$. Hence
$P(P(x))$ has at least $k$ real roots (counted without multiplicity).
$\blacksquare$
\end{solution}
\newpage

\begin{problem}{Chapter 2.2 --- Problem 151}
Let $P(x)$ be a polynomial of degree $n$. Knowing that
$$
P(k) = \frac{k}{k+1}, \quad k=0,1,\ldots,n
$$
find $P(m)$ for $m > n$.
\end{problem}

\begin{solution}{Solution}
Define $g(x) = (x+1)P(x) - x$. This is a polynomial of degree $(n+1)$
with roots $0,1,\ldots, n$. Thus for some $\lambda \in \mathbb{C}$,
$$
g(x) = \lambda x(x-1)(x-2)\cdots(x-n).
$$
Now we can re-arrange to find that
$$
P(x) = \frac{g(x) + x}{x+1}.
$$
Now since $P$ is a polynomial of degree $(n+1)$, $(x+1)$ should divide
$g(x) + x$. Factor theorem tells us that $g(-1) -1 = 0$. Now we can
compute $\lambda$,
\[
\begin{aligned}
g(-1) &= \lambda (-1)(-2)\cdots(-n-1) \\
&= \lambda (-1)^{n+1} (n+1)!
\end{aligned}
\]
Since $g(-1) = 1$ we find that $\lambda = \frac{1}{(-1)^{n+1} (n+1)!}$.
So we find that
$$
P(x) = \frac{\frac{1}{(-1)^{n+1} (n+1)!}x(x-1)\cdots(x-n) + x}{x + 1}
$$
or using generalised binomial coefficients
$$
P(x) = \frac{1}{x+1}\left( \binom{-x}{n+1} + x \right). \ \blacksquare
$$
\end{solution}
\newpage

\begin{problem}{Chapter 2.2 --- Problem 148}
Determine all polynomials $P(x)$ with real coefficients for which there
exists a positive integer $n$ such that for all $x$,
$$
P\left( x + \frac{1}{n} \right) + P\left( x - \frac{1}{n} \right) = 2P(x).
$$
\end{problem}

\begin{solution}{Solution}
The answer is that all linear polynomials $P(x) = ax + b$ work,
$a,b\in \mathbb{R}$. We verify
\[
\begin{aligned}
\left[ a\left( x + \frac{1}{n} \right) + b \right] +
\left[ a\left( x - \frac{1}{n} \right) + b \right] &=
2(ax + b).
\end{aligned}
\]
Now let's show that polynomials of degree $\geq 2$ don't work. \\

\textit{Claim 1:
Consider some} $p(x) \in \mathbb{R}[x]$, \textit{with degree} $n$.
\textit{Fix} $d > 0$.
\textit{Then} $\deg[p(x+d) - p(x)] = n-1$. \vspace{0.2cm}

\textit{Proof:}
Let
$p(x) = a_n x^n + a_{n-1} x^{n-1} + \cdots + a_0$, where $a_n \neq 0$.
Note that
\[
\begin{aligned}
p(x+d) - p(x) &= \sum_{r=0}^n a_r \left[ (x+d)^r - x^r \right] \\
&= \sum_{r=0}^n a_r \sum_{k=0}^{r-1} \binom{r}{k} x^k d^{r-k}
\end{aligned}
\]
and so $\deg [p(x+d) - p(x)] = n-1$. $\square$
\\

Suppose that $\deg P = N$.
We may re-arrange the given equation to get
$$
P\left( x + \frac{1}{n} \right) - P(x) =
P(x) - P\left( x - \frac{1}{n} \right)
$$
Define $\Delta_n(x) = P\left( x + \frac{1}{n} \right) - P(x)$. Observe
$\deg \Delta_n = N-1$. We have
$$
\Delta_n(x) - \Delta_{n-1}(x) = 0
$$
Noting that $\deg [\Delta_n(x) - \Delta_{n-1}(x)] = N-2 = 0$, we find that
$N = 2$. So now suppose that $P(x) = ax^2 + bx + c$. Plugging this into
the given equation We have
$$
\frac{2a}{n^2} = 0
$$
and so $a=0$. Thus only linear $P$ work. $\blacksquare$
\end{solution}

\end{document}
